% experimental-setup.tex -- EVM Executor Paper (BNR-2026-008)

\section{Experimental Setup}

\subsection{Implementation}

The experimental implementation consists of three Go source files targeting
go-ethereum v1.14.12:

\begin{itemize}
  \item \texttt{main.go}: Benchmark suite implementing simple transfer and
    storage write benchmarks with configurable tracing modes (no tracing,
    selective ZK tracing, full opcode capture).
  \item \texttt{opcode\_analysis.go}: Complete mapping of 75 Cancun EVM opcodes
    (plus PUSH1--32, DUP1--16, SWAP1--16) to ZK constraint estimates and
    circuit treatment recommendations.
  \item \texttt{setup.sh}: Build and execution script with dependency analysis.
\end{itemize}

\subsection{Chain Configuration}

The benchmark uses the Basis Network L2 chain configuration:

\begin{lstlisting}[language=Go,caption={Basis L2 chain configuration}]
ChainConfig{
  ChainID:      43199,  // Basis Network
  CancunTime:   0,      // Cancun from genesis
  // No PragueTime (Avalanche constraint)
}
\end{lstlisting}

The EVM target is Cancun, consistent with the Avalanche Subnet-EVM constraint
that prevents use of Pectra-era opcodes. Gas price is set to zero, matching
the zero-fee enterprise model.

\subsection{Benchmark Configurations}

Five benchmark configurations are evaluated:

\begin{enumerate}
  \item \textbf{Simple transfer, no tracing}: Baseline ETH transfers (21{,}000 gas)
    between 100 pre-funded accounts. 10{,}000 iterations with 10\% warm-up.
  \item \textbf{Simple transfer, ZK tracing}: Same workload with selective ZK
    tracer (state changes only, no opcode log).
  \item \textbf{Storage write, no tracing}: SSTORE operations via minimal
    contract (8~bytes bytecode) with varying slot addresses.
  \item \textbf{Storage write, ZK tracing}: Same workload with selective ZK
    tracer capturing StorageChange events.
  \item \textbf{Full opcode trace}: Storage write workload with complete opcode
    capture (every instruction logged).
\end{enumerate}

\subsection{StateDB Configuration}

All benchmarks use an in-memory StateDB (\texttt{rawdb.NewMemoryDatabase()})
with 100 accounts pre-funded with 1{,}000~ETH each. This eliminates disk I/O
as a confounding variable and represents the expected deployment configuration
where the L2 node maintains state in memory with periodic persistence.

\subsection{Metrics}

\begin{itemize}
  \item \textbf{Transactions per second} (tx/s): Primary throughput metric.
  \item \textbf{Average transaction time} ($\mu$s/tx): Per-transaction latency.
  \item \textbf{Trace size} (bytes/tx): Serialized JSON trace per transaction.
  \item \textbf{Tracing overhead} (\%): Throughput reduction attributable to tracing.
  \item \textbf{Memory usage} (MB): Heap allocation delta during benchmark.
\end{itemize}
